% !PRINTED 783
\weberpage{783}
%<<CONT
(in which the imaginary part of the logarithm is to be taken between $\pm \pi i$), then, if the imaginary part of $\log Q(s)$ is suitably determined, it follows that

\weq{21}{\log Q(s) = \Sigma \frac{\chi(p)}{p^s} + \frac{1}{2}\Sigma \frac{\chi(p^2)}{p^{2s}} + \frac{1}{3}\Sigma \frac{\chi(p^3)}{p^{3s}} + \cdots}

If now $A$ denotes any one of the classes (2) and $a$ a number from $A$, then by \S.\,11, (6) of this volume

\weq{22}{\overset{\chi}{\Sigma} \chi(A) \chi(n) = \Sigma \chi(an) = \mu \text{ or } = 0,}

according as $an \equiv 1$ or not $\equiv 1$ to the modulus $m$, that is, according as $n$ is contained in the class $A^{-1}$ or is not contained in it. If, then, we multiply formula (21) by $\chi(A)$ and sum with respect to $\chi$, it follows that

\weq{23}{\frac{1}{\mu} \overset{\chi}{\Sigma} \chi(A) \log Q(s) = \Sigma \frac{1}{p^s} + \frac{1}{2}\Sigma \frac{1}{p^{2s}} + \frac{1}{3}\Sigma \frac{1}{p^{3s}} + \cdots,}

where on the right the first, second, third, \ldots\ sum refers to all primes $p$ whose first, second, third, \ldots\ power is contained in $A^{-1}$.

The second part of this sum
\[
R(s) = \frac{1}{2}\Sigma \frac{1}{p^{2s}} + \frac{1}{3}\Sigma \frac{1}{p^{3s}} + \cdots
\]
is increased if each of its partial sums is extended to \gesperrt{all} primes $p$, and accordingly
\[
R(s) < \frac{1}{2}\Sigma \left(\frac{1}{p^{2s}} + \frac{1}{p^{3s}} + \cdots\right) = \frac{1}{2}\Sigma \frac{1}{p^s (p^s - 1)},
\]
and so all the more:
\[
R(s) < \frac{1}{2} \sum_{2,\infty}^{n} \frac{1}{n^{2s} (1 - n^{-s})},
\]
or, since $1 - n^{-s} > \frac{1}{2}$:
\[
R(s) < \sum_{2,\infty}^{n} \frac{1}{n^{2s}},
\]
and from this one concludes that $R$ remains finite for $s = 1$ (\S.\,191, 2.).

If we now, in formula (23),
\[
\frac{1}{\mu} \overset{\chi}{\Sigma} \chi(A) \log Q(s) = \Sigma \frac{1}{p^s} + R(s),
\]
%CONT>>
